\section{Die Nordwind-Datenbank}

\subsection{Kurz-Beschreibung}

Die Nordwind-Datenbank wird als Beispiel-Datenbank zu Microsoft Access mitgeliefert. An diesem Beispiel kann man erkennen, wie das Rechnungswesen einer Firma mit Hilfe einer Datenbank durchgeführt werden kann.

Folgendes wird durch die Tabellen der Nordwind-Datenbank abgebildet:
\begin{itemize}
\item Lieferanten liefern Artikel, die bestimmten Kategorien zugeordnet sind;
\item Kunden bestellen gewisse Artikel;
\item Die Bestellungen werden vom Personal der Firma bearbeitet;
\item Die Auslieferung einer Bestellung erfolgt über eine Versandfirma.
\end{itemize}


Die Beziehungen zwischen den Tabellen lassen sich kurz wie folgt charakterisieren:

\begin{itemize}
\item Eine Firma kann mehrere Artikel liefern, ein Artikel kommt aber von genau einer Firma
\item Eine Kategorie umfasst mehrere Artikel, jeder Artikel gehört genau einer Kategorie an
\item Ein Artikel kann in mehreren Bestellungen vorkommen und zu einer Bestellung gehören ein oder mehrere Artikel
\item Ein Kunde kann mehrere Bestellungen aufgeben, aber eine Bestellung gehört zu genau einem Kunden
\item Ein Personalmitarbeiter bearbeitet mehrere Bestellungen, aber eine Bestellung wird immer von genau einem Mitarbeiter bearbeitet
\item Eine Versandfirma liefert mehrere Bestellungen aus, aber eine Bestellung wird immer von einer Versandfirma ausgeliefert
\end{itemize}

\subsection{Tabellen}

\begin{table}[H]
\centering
\begin{tabular}{l|p{12cm}}
\textbf{Tabelle} & \textbf{Spalten} \\\midrule\midrule
Lieferant&	\emoji{key} LieferantenNr, Firma, Kontaktperson, Position, Straße, Ort, Region, PLZ, Land, Telefon, Telefax, Homepage\\\midrule
Kunde &	\emoji{key} KundenCode, Firma, Kontaktperson, Position, Straße, Ort, Region, PLZ, Land, Telefon, Telefax\\\midrule
Versandfirma &	\emoji{key} FirmenNr, Firma, Telefon\\\midrule
Personal &	\emoji{key} PersonalNr, Nachname, Vorname, Position, Anrede, Geburtsdatum, Einstellung, Straße, Ort, Region, PLZ, Land, TelefonPrivat, DurchwahlBüro, Bemerkungen, Vorgesetzter\\\midrule
Kategorie 	&\emoji{key} KategorieNr, Kategoriename, Beschreibung\\\midrule
Artikel 	&\emoji{key} ArtikelNr, Artikelname, LieferantenNr \emoji{up-right-arrow}, KategorieNr \emoji{up-right-arrow}, Liefereinheit, Einzelpreis, Lagerbestand, BestellteEinheiten, Mindestbestand, Auslaufartikel\\\midrule
Bestellung 	& \emoji{key} BestellNr, KundenCode \emoji{up-right-arrow}, PersonalNr \emoji{up-right-arrow}, Bestelldatum, Lieferdatum, Versanddatum, FirmenNr \emoji{up-right-arrow}, Frachtkosten, Empfänger, Straße, Ort, Region, PLZ Bestimmungsland\\\midrule
Bestelldetails &	BestellNr \emoji{up-right-arrow}, ArtikelNr \emoji{up-right-arrow}, Einzelpreis, Anzahl, Rabatt
\end{tabular}
\end{table}

% Stil für alle Tabellen
\newtcolorbox{dbtable}[1]{colback=blue!5!white,colframe=blue!75!black,fonttitle=\bfseries,title=#1, width=5cm}


\begin{figure}[H]
\centering
\begin{tikzpicture}
    % Tabelle: Lieferant
    \node[right] (Lieferant) {
        \begin{dbtable}{Lieferant}
            LieferantenNr \emoji{key}\\
            Firma\\
            Kontaktperson\\
            Position\\
            Straße\\
            Ort\\
            Land\\
        \end{dbtable}
    };

    % Tabelle: Kategorie
    \node[right=3cm of Lieferant] (Kategorie) {
        \begin{dbtable}{Kategorie}
            KategorieNr \emoji{key}\\
            Kategoriename\\
            Beschreibung\\
        \end{dbtable}
    };

    % Tabelle: Artikel
    \node[below=1.5cm of Lieferant] (Artikel) {
        \begin{dbtable}{Artikel}
            ArtikelNr \emoji{key}\\
            Artikelname\\
            LieferantenNr \emoji{up-right-arrow}\\
            KategorieNr \emoji{up-right-arrow}\\
            Liefereinheit\\
            Einzelpreis\\
        \end{dbtable}
    };

    % Tabelle: Kunde
    \node[right=3cm of Artikel] (Kunde) {
        \begin{dbtable}{Kunde}
            KundenCode \emoji{key}\\
            Firma\\
            Kontaktperson\\
            Position\\
            Straße\\
            Ort\\
            Land\\
        \end{dbtable}
    };

    % Tabelle: Bestellung
    \node[below=1.5cm of Artikel] (Bestellung) {
        \begin{dbtable}{Bestellung}
            BestellNr \emoji{key}\\
            KundenCode \emoji{up-right-arrow}\\
            PersonalNr \emoji{up-right-arrow}\\
            Bestelldatum\\
            Lieferdatum\\
            Versanddatum\\
            FirmenNr \emoji{up-right-arrow}\\
        \end{dbtable}
    };

    % Tabelle: Bestelldetails
    \node[right=3cm of Bestellung] (Bestelldetails) {
        \begin{dbtable}{Bestelldetails}
            BestellNr \emoji{up-right-arrow}\\
            ArtikelNr \emoji{up-right-arrow}\\
            Einzelpreis\\
            Anzahl\\
            Rabatt\\
        \end{dbtable}
    };

    % Tabelle: Personal
    \node[below=1.5cm of Bestellung] (Personal) {
        \begin{dbtable}{Personal}
            PersonalNr \emoji{key}\\
            Nachname\\
            Vorname\\
            Position\\
            Geburtsdatum\\
            Einstellung\\
            Straße\\
            Ort\\
            Land\\
        \end{dbtable}
    };

    % Tabelle: Versandfirma
    \node[right=3cm of Personal] (Versandfirma) {
        \begin{dbtable}{Versandfirma}
            FirmenNr \emoji{key}\\
            Firma\\
            Telefon\\
        \end{dbtable}
    };

    % Beziehungen
    \draw[->] (Artikel) -- (Lieferant);
    \draw[->] (Artikel) -- (Kategorie);
    \draw[->] (Bestellung) -- (Kunde);
    \draw[->] (Bestellung) -- (Personal);
    \draw[->] (Bestellung) -- (Versandfirma);
    \draw[->] (Bestelldetails) -- (Bestellung);
    \draw[->] (Bestelldetails) -- (Artikel);

\end{tikzpicture}
\end{figure}
